%% !TEX program = xelatex
\documentclass[12pt,final]{novelette} % font size, twoside, twocolumn, draft-final, use extarticle for 8-20pt

\title{Calculus III Final 2541}
\author{Kaka Cring}
\date{\today}

\begin{document}

\makeheading{ข้อสอบภาคต้น วิชา 2301217 Calculus III ปีการศึกษา 2541}

\bigskip

\begin{question}
    \item {}
    \begin{question}
        \item
        จากสมการพื้นผิว $x^2 + y^2 + z^2 + 2yz = 0$
        จงหา เมตริกซ์ของการเปลี่ยนทิศทางของแกนพิกัดที่ทำให้สมการใหม่ที่ได้อยู่ในรูป \
        $\lambda_1(x')^2 + \lambda_2(y')^2 + \lambda_3(z')^2 = k$ \\
        เมื่อเทียบกับพิกัดใหม่ $(X'Y'Z')$

        % ----------------------------------------

        % ----------------------------------------

        \item
        จงเขียนสมการใหม่ของสมการพื้นผิว $x^2 + y^2 + z^2 + 2yz + z - y = 0$
        พร้อมทั้งวาดภาพของพื้นผิวนี้ด้วย \qmark{14}

        % ----------------------------------------

        % ----------------------------------------
    \end{question}

    \item {}
    \begin{question}
        \item
        ให้ $C$ เป็นเส้นโค้ง กำหนดโดยสมการ
        $$
            \vec{r}(t) = \left(\frac{t^3}{3}, t^2, 2t\right), \quad t \in [0, 3]
        $$
        จงหาความยาวของเส้นโค้ง $C$ และความโค้งของเส้นโค้ง $C$ ณ จุด $\vec{r}(t)$ ใด ๆ \qmark{8}

        % ----------------------------------------

        % ----------------------------------------

        \item
        ให้ $C$ เป็นเส้นโค้ง กำหนดโดยสมการ
        $$
            \vec{r}(t) = \left(\frac23 t^3, \frac{t^2}{2}, 2t\right), \qquad t \in [0, 3]
        $$
        จงหาเวกเตอร์หนึ่งหน่วยสัมผัส เวกเตอร์นอร์มัลหลัก และเวกเตอร์นอร์มัลคู่ ณ จุด $\vec{r}(1)$ \qmark{8}

        % ----------------------------------------

        % ----------------------------------------
    \end{question}

    \item {}
    \begin{question}
        \item
        จงเขียน $\iiint_D xz^2 \dd{v}$ ในระบบพิกัดทรงกลม เมื่อ $D$ เป็นวัตถุทรงตัน ที่อยู่ภายในทรงกลม \\
        $x^2 + y^2 + (z-2)^2 = 4$
        และอยู่ภายนอกกรวย $\sqrt3 z = \sqrt{x^2 + y^2}$ \\
        (เขียนภาพโดเมนของการอินทิเกรตและไม่ต้องคำนวณค่า) \qmark{6}

        % ----------------------------------------

        % ----------------------------------------

        \item
        ให้ $S$ เป็นทรงตันที่ปิดล้อมด้วยผิว $y = 4x^2 + 9z^2$, $z = x^2$, $z = 1$ และ $y = 0$ \\
        จงเขียน $\iiint_S f(x, y, z) \dd{v}$ โดยใช้อินทิกรัลซ้อน โดยมีลำดับการอินทิเกรต $\dd{y} \dd{z} \dd{x}$ \qmark{6}

        % ----------------------------------------

        % ----------------------------------------

        \item
        จงเขียนอินทิกรัล \quad $\displaystyle \int\limits_0^{\sqrt3} \int\limits_{-\sqrt{3-y^2}}^{\sqrt{3-y^2}} \int\limits_{x^2+y^2}^{2+\sqrt{4-x^2-y^2}} f(x, y, z) \dd{z}\dd{x}\dd{y}$ \\[.5em]
        ในระบบพิกัดทรงกระบอก \qmark{6}

        % ----------------------------------------

        % ----------------------------------------
    \end{question}

    \item {}
    จงเปลี่ยนอินทิกรัล $\displaystyle \int\limits_{-2}^2 \int\limits_{-\sqrt{4-x^2}}^0 \int\limits_0^{\sqrt{4-x^2-y^2}} z^2 \sqrt{x^2 + y^2 + z^2} \dd{z} \dd{y} \dd{x}$ \\[.5em]
    ให้อยู่ในระบบพิกัดทรงกลม พร้อมทั้งเขียนโดเมนของการอินทิเกรต และหาค่าอินทิกรัลด้วย \qmark{8}

    % ----------------------------------------

    % ----------------------------------------

    \item {}
    \begin{question}
        \item
        จงหาค่า $\int_C (x + y + z - 1) \dd{s}$ เมื่อ $C$ เป็นส่วนโค้ง $y = x^2, z + y = 1$ จากจุด $(0,0,1)$ ไปยังจุด $(1,1,0)$ \qmark{8}

        % ----------------------------------------

        % ----------------------------------------

        \item
        จงหาค่า $\displaystyle \int_C \frac{x}{x^2+y^2+1} \dd{x} + \frac{y}{x^2+y^2+1} \dd{y}$ เมื่อ $y$ เป็นเส้นโค้ง
        กำหนดโดยสมการ \\[.5em]
        $\displaystyle \vec{r}(t) = (e^t \cos t, e^t), \quad 0 \leq t \leq \frac{\pi}{2}$ \qmark{8}

        % ----------------------------------------

        % ----------------------------------------
    \end{question}

    \item {}
    ให้ $C$ เป็นอาณาบริเวณที่ปิดล้อมด้วยเส้นโค้ง $y = x^3$ เส้นตรง $x + y = 2$ และ $x = 0$ ในทิศทางทวนเข็มนาฬิกา

    จงหา $\int_C \vec{F} \cdot \dd{\vec{r}}$ เมื่อ $\vec{F}(x, y) = (x^3y^2 + e^{x^2}, \tan y + e^y + yx^2)$ \qmark{8}

    % ----------------------------------------

    % ----------------------------------------

    \item {}
    จงหาค่าของ $\iint_S (1-x)yz \dd{s}$ โดยให้พื้นผิว $S$ เป็นส่วนหนึ่งของพื้นผิว $y^2 = 1-x$, $x \geq 0$ และ $0 \leq z \leq 3$ \qmark{8}

    % ----------------------------------------

    % ----------------------------------------

    \item {}
    กำหนดให้ $\vec{F}(x, y, z) = (y^2 + z^2 + 6zx, 2y(x + e^{y^2}) + 3y^2z, 2xz + y^3 + 3x^2)$

    \begin{question}
        \item
        จงหา $\int_C \vec{F} \cdot \dd{\vec{r}}$ เมื่อ $C$ เป็นเส้นโค้งที่มีสมการ
        $$
            \vec{r}(t) = (t \sin t, 3 \cos t, 2 \cos t - \sin t), \quad t \in \left[0, \frac\pi2\right]
        $$

        % ----------------------------------------

        % ----------------------------------------

        \item
        จงหา $\varointctrclockwise_C \vec{F} \cdot \dd{\vec{r}}$  เมื่อ $C$ เป็นวงรีที่มีสมการ
        $$
            \frac{(x-1)^2}{25} + \frac{(y+3)^2}{16} = 1 \quad \text{และ} \quad z = 1 \quad \text{โดยอินทิเกรตในทิศทางทวนเข็มนาฬิกา}
        $$
        \qmark{12}

        % ----------------------------------------

        % ----------------------------------------
    \end{question}
\end{question}

\pagebreak

\noindent
จงใช้\uline{ทฤษฎีบทของกรีน}หาพื้นที่ของบริเวณที่ปิดล้อมด้วยเส้นโค้ง $y = x^2$ และ $y = x$
[ข้อสอบปลายภาคต้น ปีการศึกษา 2549 / 14]

% ----------------------------------------

% ----------------------------------------

\end{document}